\documentclass[tikz,border=18pt]{standalone}
\usepackage{luatexja}
\usepackage{tikz}
\usetikzlibrary{shapes.geometric,arrows.meta,positioning,calc}

\definecolor{fidoamber}{RGB}{215, 120, 20}
\definecolor{goodgreen}{RGB}{30, 135, 55}
\definecolor{clientblue}{RGB}{20, 80, 150}

\begin{document}
\begin{tikzpicture}[>=Stealth, font=\small]

  % 端末側（生体・所持・秘密鍵）
  \node[draw=clientblue, fill=blue!6, rounded corners=8pt, line width=1.5pt, minimum width=9.2cm, minimum height=7.2cm, align=center] (client_box) at (-5.4, 2.0) {};
  \node[above, font=\large\bfseries, text=clientblue] at (-5.4, 4.8) {【ユーザー端末】スマホ・PC（生体認証＋Secure Enclave）};

  \node[draw=clientblue, fill=white, rounded corners=6pt, thick, align=center, minimum width=7.6cm, minimum height=1.5cm] (bio_step) at (-5.4, 3.6) {
    \textbf{ローカル生体認証 (Touch ID / Face ID)}\\[2pt]
    \footnotesize ※生体データは端末外に一切送信されない！
  };

  \node[draw=goodgreen, fill=green!15, rounded corners=6pt, very thick, align=center, minimum width=7.6cm, minimum height=2.2cm] (key_step) at (-5.4, 0.8) {
    \textbf{\textcolor{goodgreen}{\large 端末内チップの「秘密鍵」で署名}}\\[3pt]
    \footnotesize サーバーから届いたランダムな合言葉（Challenge）に\\
    \footnotesize 端末内の秘密鍵でデジタル署名を作成！\\[2pt]
    \scriptsize $\Rightarrow$ 秘密鍵はチップ外に絶対に出ない（流出ゼロ！）
  };

  \draw[->, very thick, goodgreen] (bio_step.south) -- (key_step.north)
    node[midway, right=4pt, font=\scriptsize\bfseries, text=goodgreen] {本人確認完了};

  % サーバー側（公開鍵・チャレンジ生成・署名検証）
  \node[draw=fidoamber, fill=orange!6, rounded corners=8pt, line width=1.5pt, minimum width=9.2cm, minimum height=7.2cm, align=center] (server_box) at (5.4, 2.0) {};
  \node[above, font=\large\bfseries, text=fidoamber] at (5.4, 4.8) {【サービス提供者】Webサーバー・クラウド};

  \node[draw=fidoamber, fill=white, rounded corners=6pt, thick, align=center, minimum width=7.6cm, minimum height=1.5cm] (chal_step) at (5.4, 3.6) {
    \textbf{① 毎回異なる使い捨ての合言葉}\\[2pt]
    \footnotesize (Challenge コードを発行して端末へ送信)
  };

  \node[draw=fidoamber, fill=orange!15, rounded corners=6pt, very thick, align=center, minimum width=7.6cm, minimum height=2.2cm] (verify_step) at (5.4, 0.8) {
    \textbf{\textcolor{fidoamber}{\large ② 登録済みの「公開鍵」で署名検証}}\\[3pt]
    \footnotesize 署名が正しければログイン成功！\\[2pt]
    \scriptsize サーバーには「公開鍵」しか保存されていないため\\
    \scriptsize サーバーがハッキングされてもパスワード流出は原理的に不可能！
  };

  \draw[->, thick, dashed, clientblue] (chal_step.west) -- (bio_step.east)
    node[midway, above=4pt, font=\scriptsize\bfseries, text=clientblue] {Challenge 送信};

  \draw[->, very thick, goodgreen] (key_step.east) -- (verify_step.west)
    node[midway, above=4pt, font=\scriptsize\bfseries, text=goodgreen] {署名付き応答 (Response)};

  % 最下部：フィッシング耐性
  \node[draw=goodgreen, fill=green!25, rounded corners=6pt, line width=1.5pt, align=center, minimum width=20.0cm, inner sep=6pt] at (0, -2.4) {
    \textbf{\Large ★ なぜパスキーは「フィッシング詐欺」を100\%無効化できるのか？}\\[3pt]
    \small 端末はブラウザの正規ドメイン（FQDN）と厳格にバインドして署名するため、偽サイトにいくら騙されても署名が絶対に作成・送信されない！
  };

\end{tikzpicture}
\end{document}
